\section{Formalization architecture}\label{sec:formalization}

The Lean~4 formalization is split between the shared repository package \texttt{lean-machine}
(\texttt{../lean-machine}) and this paper's package \texttt{BoundedArityDecoders} (\texttt{lean/}),
both pinned to Lean~4 and Mathlib version \texttt{v4.34.1} \cite{Mathlib2026}. Through Lake path
dependencies in \leanfile{lean/lakefile.toml}, \texttt{BoundedArityDecoders} requires
\leanname{SharedDecoders} (\cite{SharedDecoders2026}, which transitively imports
\leanname{SwapRigidity} from \cite{SeventeenthPaper2026} and \leanname{Transpose} from
\cite{FirstPaper2026}), \leanname{ClassUniformMemo} (\cite{ClassUniformMemo2026}), and
\leanname{LeanMachine}. None of the earlier papers is modified,
and in \texttt{lean-machine} only the package root, axiom audit, and \texttt{README} are appended to
register four new reusable modules under \leanfile{LeanMachine/ShuffleShear/}:
\begin{itemize}[leftmargin=1.5em]
\item \leanfile{Model.lean}, \leanfile{StoneShuffle.lean}, \leanfile{StoneUnshuffle.lean}, and
  \leanfile{RowShear.lean}: standalone deterministic multitape Turing machine semantics and verified
  linear-time machines for Stone's shuffle ($3$ work tapes, $\le 150 L$ steps), Stone's unshuffle
  ($2$ work tapes, $\le 150 L$ steps), and within-row cyclic shear ($2$ work tapes, $\le 200 L$
  steps).
\end{itemize}

Inside \leanfile{lean/BoundedArityDecoders/}, twenty-seven public modules are organized into three layers:
\begin{enumerate}[label=(\arabic*)]
\item \emph{Nine definition and specification modules}
  (\leanfile{AffineCubeComposition.lean},
  \leanfile{SumsetRankDeficit.lean},
  \leanfile{ExtensionLinearRun.lean},
  \leanfile{ExtensionAperiodicSix.lean},
  \leanfile{ExtensionTwoSidedProfile.lean},
  \leanfile{SwapCodeAnswers.lean},
  \leanfile{OperatorDuality.lean},
  \leanfile{NativeSweepInterface.lean}, and
  \leanfile{BoundedAritySetCapacity.lean}), defining the algebraic operators, one-stage linear
  run-extension invariants, coarse--fine interval decompositions, syntactic causal cones, and named
  proposition types;
\item \emph{Eleven primary proof modules}
  (\leanfile{AffineCubeCompositionProofs.lean},
  \leanfile{SumsetRankDeficitProofs.lean},
  \leanfile{ExtensionLinearRunSingleShift.lean},
  \leanfile{ExtensionLinearRunSharpness.lean},
  \leanfile{ExtensionLinearRunProofs.lean},
  \leanfile{ExtensionAperiodicSixProofs.lean},
  \leanfile{ExtensionTwoSidedProfileProofs.lean},
  \leanfile{SwapCodeAnswersProofs.lean},
  \leanfile{OperatorDualityProofs.lean},
  \leanfile{NativeSweepInterfaceProofs.lean}, and
  \leanfile{BoundedAritySetCapacityProofs.lean}), proving the two-stage composition theorems,
  the one-stage linear run-extension and $\Z_6$ bounds, the two-sided coarse--fine swap-code bound,
  and the finite controls; and
\item \emph{Seven consequence, transfer, and bridge modules}
  (\leanfile{ExtensionFieldTransfer.lean},
  \leanfile{AffineCubeConsequences.lean},
  \leanfile{SumsetRankConsequences.lean},
  \leanfile{DualityAndSwapConsequences.lean},
  \leanfile{ShuffleShearBridge.lean},
  \leanfile{MachineSweepConsequences.lean}, and
  \leanfile{SetCapacityConsequences.lean}), proving the algebraic-closure primitive-root and
  scalar-extension transfer lemmas, the explicit constants $c(M_1, M_2)$ and
  $D_0(M_1, M_2)$, the general-field finite controls at $d = 13$, $d = 4$, and $d = 12$, the
  \leanname{TensorProduct} anti-diagonal equivalences, the transport of
  \leanname{LeanMachine.ShuffleShear} to \leanname{Transpose.TM}, the symbolic
  \lean{CausalConeTree} evaluator, the $1/(\Delta+1)$ random-order bound, and the negative
  per-column cut-absorption bounds.
\end{enumerate}

Every numbered theorem, lemma, proposition, corollary, and example in the manuscript is proved as a
\texttt{theorem} declaration in these modules; there are no unformalized or partially formalized
numbered results.
